\cleardoublepage

\chapter*{Abstract\markboth{Abstract}{Abstract}}

Security in modern systems is a constantly evolving field in which data integrity and control over system operation play a fundamental role.
This is due to the presence of threats and external agents whose malicious purpose is to compromise such information, making it necessary to develop increasingly robust protection mechanisms.\\

In this context, one of the most critical threats currently affecting information systems is represented by \textit{rootkits}, malicious software designed to hide and manipulate system resources or activity, such as processes, files or network connections.
These threats may operate both in \textit{User Space} and in \textit{Kernel Space}, compromising the information received by both the administrator and the user and preventing them from obtaining a truthful view of the real state of the system.
This undermines the integrity of monitoring and control tasks, potentially leading the administrator to trust information sources that may have been compromised.\\

To address this problem, this project designs and implements a \textit{Loadable Kernel Module} (\textit{LKM}) on a \textit{Linux} system with the aim of providing an authenticated source of critical system information against possible manipulation.
The solution introduces a cryptographic authentication mechanism based on \textit{HMAC} for the content provided by the \textit{kernel}, using a secret symmetric key shared only between the administrator and the module.
This approach relies on a chain of trust in which it is assumed that the \textit{kernel} and the \textit{LKM} have not been modified, for example through signing mechanisms for both the \textit{kernel} and the module, thereby reinforcing the integrity of the information provided.\\

The experiments carried out approximate a real operating environment in which sensitive information is requested from \textit{User Space} to \textit{Kernel Space}.
For this purpose, the sensitive system content to be authenticated is defined as a list of all running processes and their corresponding metadata, generated inside the \textit{Loadable Kernel Module} itself and used as an example of critical information that may be manipulated by a \textit{rootkit}.
An \textit{HMAC} is computed over this content using the corresponding secret key and is attached to the information provided to the user, so that its integrity and authenticity can be verified.\\

To validate this proposal, user-level components were developed to test the functionality. A first component creates a file whose descriptor is transferred to the \textit{kernel} so that it can reference it, thus establishing the secure communication channel with the core of the system.
On the other hand, the integrity of the reported data is confirmed through a second verification component that analyzes the information written by the \textit{kernel}, locally computes the \textit{HMAC} of the content using the same secret key and then compares it with the received value to ensure that both match and that the content is authentic.\\
